%! Exponential Function Manipulation — Unit 4, Lesson 4.4 — Student Packet Cover Sheet
\documentclass[10pt]{article}
\usepackage{precalculus-article}
\usepackage{precalculus-boxes}
\usepackage{ltablex}
\keepXColumns

\begin{document}

% Full-bleed navy banner
\begin{tikzpicture}[remember picture, overlay]
  \fill[navy] (current page.north west)
    rectangle ([yshift=-0.9in]current page.north east);
\end{tikzpicture}
\vspace{-0.8in}

\begin{center}
  {\color{white}\textbf{\LARGE Precalculus}}\\[4pt]
  {\color{skymid}\large\textbf{Unit 4: Exponential Functions}}\\[6pt]
  {\color{white}\textbf{\large Lesson 4.4 \quad Exponential Function Manipulation}}
\end{center}

\vspace{0.22in}
\namedateperiod
\vspace{0.18in}

\begin{learningtargetbox}
\small
\begin{enumerate}[leftmargin=1.5em, itemsep=5pt]
  \item I can apply the product property $b^m \cdot b^n = b^{m+n}$ to rewrite a horizontal
        translation $b^{x+k}$ as a vertical dilation $a \cdot b^x$.
  \item I can apply the power property $(b^m)^n = b^{mn}$ to rewrite a horizontal dilation
        $b^{cx}$ as a single base raised to the $x$ power.
  \item I can simplify exponential expressions using negative exponents and fractional
        exponents (roots).
\end{enumerate}
\end{learningtargetbox}

\vspace{0.14in}

\begin{tocbox}
\small
\renewcommand{\arraystretch}{1.5}
\begin{tabularx}{\linewidth}{@{} c l X @{}}
\rowcolor{sky}
\textbf{\#} & \textbf{Component} & \textbf{Description} \\
\midrule
1 & Warm-Up        & Spiral review: exponential form, growth/decay, table to equation \\
2 & Guided Notes   & Product property; power property; negative \& fractional exponents \\
3 & Group Activity & Rewriting exponential expressions (tiered R / A / E) \\
4 & Exit Ticket    & Rewrite using product, power, and exponent properties \\
5 & Homework       & Mixed: translations, base changes, exponent rules, AP-style MC. \textbf{Scored 1 to 5.} \\
\end{tabularx}
\end{tocbox}

\end{document}
